\section{Setting, methods and evaluation protocol}
\subsection{Setting and conformal regions}
Let $(X,Y)$ have $Y\in\simplex=\{y\in\mathbb{R}^K_{\ge0}:\sum_k y_k=1\}$, and let $\hat p:\mathcal X\to\simplex$ be a predictor trained on a fit split. Given a score $s(y,p)$ (larger = worse fit), a calibration set of size $n$ from a split independent of the fit split, and level $\alpha$, the split conformal threshold is $\hat q$, the $k$-th smallest calibration score with $k=\lceil (n+1)(1-\alpha)\rceil$ ($\hat q=\infty$ if $k>n$), and the region is $C(x)=\{y\in\simplex: s(y,\hat p(x))\le\hat q\}$.

\paragraph{Scores.} With $\tilde y=\mathrm{closure}(y+\varepsilon)$ for zero replacement: Euclidean $\|y-p\|_2$; total variation $\tfrac12\|y-p\|_1$; Hellinger $\|\sqrt y-\sqrt p\|_2$; $\mathrm{KL}(y\|p)$; $\|y-p\|_\infty$; Aitchison $\|\clr(\tilde y)-\clr(\tilde p)\|_2$; and the Dirichlet negative log-density $-\log \mathrm{Dir}(\tilde y;\hat\phi \tilde p)$ with a single precision $\hat\phi$ fitted by maximum likelihood on a separate split.

\paragraph{Calibration schemes.} \emph{Plain} split conformal; \emph{normalised} scores $s/\hat\sigma(x)$ where $\hat\sigma$ is a random forest for the score fitted on a third split; \emph{Mondrian} calibration with one threshold per predicted-dominant-part group; and \emph{CV+} with $10$ folds for the Euclidean score and a multinomial-logit predictor.

\subsection{Evaluation protocol}\label{sec:protocol}
\paragraph{Marginal validity.} For each configuration we repeat the whole pipeline (fresh data draw, fit, calibrate, test) $R$ times and report the mean test coverage with a confidence interval over replicates, and compare the empirical distribution of calibration-conditional coverage with the Beta law below.
\paragraph{Size.} A region's size is its Lebesgue measure on the simplex relative to the measure of the simplex, estimated by uniform Monte Carlo draws on $\simplex$ for $K\le5$ only; for larger $K$ the hit fraction vanishes and no volume is reported. Sizes under different reference measures (for example an Aitchison ball) are not compared with each other.
\paragraph{Conditional coverage.} Coverage is estimated within strata of the predicted maximal part and of the predicted entropy (quantile bins), within a boundary stratum ($\min_k y_k<0.01$), by worst-slab coverage \citep[a slab chosen on one half of the test data and evaluated on the other;][]{romano2020classification}, and, in the heteroscedastic worlds, by the true noise scale. On synthetic data these are measured against a known conditional law. We report the range of bin coverages (min/max), and calibrate their interpretability by a null simulation (Section~\ref{sec:exp}).

\subsection{Theory}\label{sec:theory}
\begin{proposition}[Coverage law of split conformal prediction]\label{prop:beta}
Let the calibration scores $S_1,\dots,S_n$ and the test score $S_{n+1}$ be exchangeable with continuous common distribution function $F$, independent of the fit split, and let $\hat q=S_{(k)}$ with $k=\lceil (n+1)(1-\alpha)\rceil\le n$. Then the coverage conditional on the calibration set, $V=\Pr(S_{n+1}\le\hat q\mid S_{1:n})=F(S_{(k)})$, has the law $\mathrm{Beta}(k,\,n+1-k)$. In particular $\mathbb E V=k/(n+1)\in[1-\alpha,\,1-\alpha+\tfrac1{n+1})$ and the law does not depend on $F$, hence not on the score, the predictor or the response family.
\end{proposition}
\begin{proof}
$F(S_i)$ are i.i.d.\ uniform by the probability integral transform, and the $k$-th order statistic of $n$ i.i.d.\ uniforms is $\mathrm{Beta}(k,n+1-k)$. The mean follows from $k\ge(n+1)(1-\alpha)$ and $k<(n+1)(1-\alpha)+1$.
\end{proof}
This is the finite-sample statement used to check the pipeline \citep{angelopoulos2021gentle}. It requires continuity of the score distribution: with exact zeros or rounded parts, ties can occur and $V$ stochastically dominates the Beta law (coverage $\ge1-\alpha$ but possibly conservative). Group-wise (Mondrian) calibration obeys the proposition within each group, with $n$ replaced by the random group size. Normalised scores stay valid because $\hat\sigma$ is fitted on a separate split. CV+ carries only a weaker guarantee ($\ge 1-2\alpha$ in general) and is not covered by the proposition.

\paragraph{Invariances.} Sets are unchanged under any strictly increasing transformation of the score (for instance Euclidean distance and its square), and the Aitchison ball is the same in every orthonormal ILR basis, since the Euclidean distance of CLR coordinates equals that of any orthonormal coordinates in the $(K-1)$-dimensional hyperplane. Only the \emph{shape} of the score matters for size and conditional coverage, not its monotone reparametrisation.


\begin{proposition}[Vertex labels]\label{prop:vertex}
If the response is a vertex $e_j$ (a class label), then total variation, Hellinger and $\mathrm{KL}(e_j\|p)$ are strictly decreasing functions of $p_j$ alone ($\tfrac12\|e_j-p\|_1=1-p_j$, $\|\sqrt{e_j}-\sqrt p\|_2^2=2-2\sqrt{p_j}$, $\mathrm{KL}=-\log p_j$), so the sets they induce coincide with those of the LAC score $1-p_j$. The Euclidean (Brier) score $\|e_j-p\|_2^2=1-2p_j+\|p\|_2^2$ and the Aitchison score depend on the rest of $p$, and generally give different sets.
\end{proposition}
This identifies which simplex scores are genuinely different from LAC on classifier outputs.

\paragraph{Pointwise conditional coverage under Dirichlet noise (large precision).} If $Y\mid x\sim\mathrm{Dir}(\phi\,\mu(x))$ and the predictor is $\mu(x)$, then for large $\phi$ the residual is approximately Gaussian with covariance $\Sigma_{\mathrm{L2}}=(\mathrm{diag}\,\mu-\mu\mu^\top)/(\phi+1)$ in raw coordinates, $\Sigma_{\mathrm{Ait}}=G\,\mathrm{diag}(\mu)^{-1}G/(\phi+1)$ in CLR coordinates ($G=I-\mathbf 1\mathbf 1^\top/K$), and $\Sigma_{\mathrm{Hel}}=(I-\sqrt\mu\sqrt\mu^\top)/(4(\phi+1))$ in Hellinger coordinates. Writing $\lambda_j(x)$ for the eigenvalues of $\Sigma(x)$, the squared score is approximately $\sum_j\lambda_j(x)Z_j^2$ with $Z_j$ i.i.d.\ standard normal; the marginal threshold $q$ is the $(1-\alpha)$ quantile of the $x$-mixture of these laws, and the predicted conditional coverage at $x$ is $\Pr(\sum_j\lambda_j(x)Z_j^2\le q^2)$. We check this first-order approximation by simulation in Section~\ref{sec:exp} and report where it fails; it is an approximation, not a theorem.
